\chapter{Operações Lógicas}

\begin{questao}{Questão 7}

Qual o valor lógico da proposição ``Não é verdade que 12 é um número ímpar''?

\begin{alternativas}
    \item Verdadeiro.
    \item Falso.
\end{alternativas}

\end{questao}

\begin{questao}{Questão 8}

Sejam $p$ e $q$ duas proposições. Se $V(p \wedge q) = V$, logo podemos afirmar que

\begin{alternativas}
    \item $V(p) = F$
    \item $V(p) = V$
    \item $V(q) = F$
    \item Não é possível definir a valoração das proposições.
\end{alternativas}

\end{questao}

\begin{questao}{Questão 9}

Sejam as proposições $p$: ``Está frio'' e $q$: ``Está chovendo''. Em linguagem natural, $\sim p \wedge \sim q$ é traduzida como

\begin{alternativas}
    \item Não está calor ou nem está chovendo.
    \item Não está frio ou está chovendo.
    \item Não está calor e nem está chovendo.
    \item Não está frio e nem está chovendo.
\end{alternativas}

\end{questao}

\begin{questao}{Questão 10}

Sejam as proposições $p$: ``Jorge é rico'' e $q$: ``Carlos é feliz''. Em linguagem simbólica, ``Jorge não é rico, mas Carlos é feliz'' é expresso como

\begin{alternativas}
    \item $\sim p \vee q$
    \item $\sim p \vee \sim q$
    \item $\sim p \wedge q$
    \item $\sim p \wedge \sim q$
\end{alternativas}

\end{questao}

\begin{questao}{Questão 11}

Sejam $p$ e $q$ duas proposições. Se $V(\sim p \vee q) = V$ e $V(p) = F$, logo podemos afirmar que

\begin{alternativas}
    \item $V(q) = F$
    \item $V(q) = V$
    \item $V(p) = V$
    \item Não é possível definir $V(q)$
\end{alternativas}

\end{questao}

\begin{questao}{Questão 12}

Sejam $p$ e $q$ duas proposições. Se $V(p \vee q) = V$ e $V(p \wedge q) = F$, logo podemos afirmar que

\begin{alternativas}
    \item Pelo menos uma das duas proposições é falsa.
    \item $V(p) = V$
    \item $V(q) = V$
    \item Não é possível afirmar nada sobre as proposições.
\end{alternativas}

\end{questao}